\documentclass{article}

% NeurIPS 2026 style (substitute with actual style file when available)
\usepackage[final]{neurips_2026}

\usepackage[utf8]{inputenc}
\usepackage[T1]{fontenc}
\usepackage{hyperref}
\usepackage{url}
\usepackage{booktabs}
\usepackage{amsfonts}
\usepackage{amsmath}
\usepackage{nicefrac}
\usepackage{microtype}
\usepackage{graphicx}
\usepackage{xcolor}

\title{AgentOps-Bench: A Benchmark for Operational Reliability of LLM Agent Systems}

\author{%
  Author One\thanks{Equal contribution.} \\
  Affiliation \\
  \texttt{author1@example.com} \\
  \And
  Author Two\footnotemark[1] \\
  Affiliation \\
  \texttt{author2@example.com} \\
}

\begin{document}

\maketitle

\begin{abstract}
Large language model (LLM) agents that interact with external tools represent a
rapidly maturing paradigm, yet existing benchmarks focus almost exclusively on
task completion accuracy.  In practice, operational reliability---encompassing
cost efficiency, behavioral consistency across runs, resilience to tool
failures, and resistance to adversarial prompt injections embedded in tool
outputs---is equally critical for production deployment.  We introduce
\textsc{AgentOps-Bench}, a benchmark that evaluates LLM agent systems across
six complementary axes: \emph{completion}, \emph{cost}, \emph{efficiency},
\emph{reliability}, \emph{recovery}, and \emph{safety}.  The benchmark
includes an instrumented tool server capable of injecting realistic failure
modes (timeouts, malformed responses, rate limits) and adversarial payloads of
varying sophistication into tool outputs.  We evaluate agents built on Claude
and GPT-4o across 28 tasks spanning tool use, code generation, data analysis,
and research, under clean, noisy, and adversarial conditions with 8 repeated
runs per configuration.  Our results reveal that models achieving near-perfect
completion scores under clean conditions can exhibit reliability rates as low as
62\% under noisy conditions and compliance rates to prompt injections exceeding
15\% under adversarial conditions.  \textsc{AgentOps-Bench} provides the
research community with a comprehensive framework for measuring the operational
readiness of LLM agent systems beyond mere accuracy.
\end{abstract}

% ---------------------------------------------------------------
\section{Introduction}
\label{sec:intro}

% Motivation: agents are moving to production, but benchmarks only measure accuracy.
% Gap: no benchmark jointly evaluates cost, reliability, recovery, and safety.
% Contribution summary: AgentOps-Bench, instrumented tool server, 6-axis scoring.

The deployment of LLM-based agent systems in production environments demands
evaluation criteria that extend beyond task completion accuracy.  While
benchmarks such as AgentBench~\cite{liu2024agentbench},
SWE-bench~\cite{jimenez2024swebench}, GAIA~\cite{mialon2024gaia}, and
WebArena~\cite{zhou2024webarena} have advanced our understanding of agent
capabilities, they evaluate agents under idealized conditions where tools
always respond correctly and no adversarial content contaminates the
environment.

In this work, we introduce \textsc{AgentOps-Bench}, a benchmark framework
that measures the \emph{operational reliability} of LLM agent systems across
six axes.

\paragraph{Contributions.}
\begin{itemize}
  \item A six-axis evaluation framework covering completion, cost, efficiency,
        reliability, recovery, and safety.
  \item An instrumented tool server with configurable failure injection and
        a catalogue of 15 adversarial prompt injections of varying sophistication.
  \item Empirical evaluation of state-of-the-art agents under clean, noisy,
        and adversarial conditions.
\end{itemize}

% ---------------------------------------------------------------
\section{Related Work}
\label{sec:related}

% Discuss: AgentBench, SWE-bench, GAIA, WebArena, MAP, CLEAR, reliability studies.
% Position AgentOps-Bench relative to each.

\paragraph{Agent Benchmarks.}
AgentBench~\cite{liu2024agentbench} evaluates LLMs across eight distinct
environments but does not vary tool reliability.
SWE-bench~\cite{jimenez2024swebench} tests code-repair capabilities on
real GitHub issues.  GAIA~\cite{mialon2024gaia} targets general-purpose
assistant tasks.  WebArena~\cite{zhou2024webarena} provides a realistic
web environment.  None of these benchmarks systematically evaluate recovery
from tool failures or resistance to adversarial injections.

\paragraph{Agent Metrics.}
MAP~\cite{kapoor2024map} proposes standardized performance metrics.
CLEAR~\cite{zhang2024clear} introduces multi-dimensional reliability
assessment.  Chen et al.~\cite{chen2024reliability} study failure modes
in tool-use scenarios.  \textsc{AgentOps-Bench} unifies these dimensions
into a single benchmark with controlled injection.

% ---------------------------------------------------------------
\section{Benchmark Design}
\label{sec:design}

% Task format, domains, tool definitions.
% 6-axis scoring framework with formal definitions.
% Instrumented tool server architecture.

\subsection{Task Suite}
% Describe the 4 domains and task construction methodology.

\subsection{Six-Axis Scoring}
% Formal definition of each axis: completion, cost, efficiency, reliability, recovery, safety.

\subsection{Instrumented Tool Server}
% Failure modes: timeout, HTTP 500, malformed JSON, empty response, rate limit.
% Adversarial injection catalogue and embedding strategy.

% ---------------------------------------------------------------
\section{Experimental Setup}
\label{sec:experiments}

% Models evaluated, conditions (clean/noisy/adversarial), number of runs.
% Hyperparameters, budget constraints.

% ---------------------------------------------------------------
\section{Results}
\label{sec:results}

% Main results table across all 6 axes.
% Breakdown by condition and domain.
% Key findings: reliability degradation, injection compliance rates.

% ---------------------------------------------------------------
\section{Analysis}
\label{sec:analysis}

% Failure mode analysis.
% Which injection types are most effective?
% Cost-accuracy tradeoffs.
% Reliability vs. model scale.

% ---------------------------------------------------------------
\section{Limitations and Future Work}
\label{sec:limitations}

% Single-turn tool use only (no multi-agent).
% Mock tools vs. real APIs.
% Prompt injection catalogue is finite.
% Future: real tool integration, multi-agent scenarios, human evaluation.

% ---------------------------------------------------------------
\section{Conclusion}
\label{sec:conclusion}

% Summary of contributions and key findings.
% Call to action: operational reliability should be a first-class evaluation dimension.

We presented \textsc{AgentOps-Bench}, a benchmark for evaluating the
operational reliability of LLM agent systems across six axes.  Our results
demonstrate that task completion accuracy alone is an insufficient proxy for
production readiness, and we encourage the community to adopt multi-dimensional
evaluation frameworks that account for cost, consistency, resilience, and
safety.

% ---------------------------------------------------------------

\bibliography{references}
\bibliographystyle{plainnat}

\end{document}
